% !PRINTED 582
\weberpage{582}

Nun ist für variable $\varkappa$ und $v$ [\S\,45, (5)]:
\weq{8}{\mathrm{dn}\,(v,\varkappa) = \dfrac{\mathrm{dn}\,(iv,\varkappa')}{\mathrm{cn}\,(iv,\varkappa')}.}
Setzen wir hierin:
\[
v = \Omega_{h,h'} = \dfrac{2hK + 2h'iK'}{m},
\]
\[
iv = \dfrac{-2h'K' + 2hiK}{m} = \Omega'_{-h',h},
\]
so haben die $\Omega'$ dieselbe Bedeutung für $\varkappa'$ wie die $\Omega$ für $\varkappa$, d.\,h. sie entsprechen der Vertauschung von $(\omega, -1/\omega)$ und es ergiebt sich aus (7) und (8), da\ss{} auch
\[
\dfrac{\mathrm{dn}\,\Omega_{h,h'}}{\sqrt{\varkappa'}} \text{ eine Einheit ist.}
\]
Es sind also auch
\weq{9}{\begin{aligned}
\dfrac{1}{\sqrt{\varkappa'}}\,\mathrm{dn}\,\Omega_{h,h'}
&= \dfrac{\vartheta_{00}\left(\dfrac{h+h'\omega}{m}\right)}{\vartheta_{01}\left(\dfrac{h+h'\omega}{m}\right)}, \\
\sqrt{\dfrac{\varkappa}{\varkappa'}}\,\mathrm{cn}\,\Omega_{h,h'}
&= \dfrac{\vartheta_{10}\left(\dfrac{h+h'\omega}{m}\right)}{\vartheta_{01}\left(\dfrac{h+h'\omega}{m}\right)}\ \text{algebraische Einheiten.}
\end{aligned}}
Aus der Formel \S\,155, (20) folgt für $v=0$:
\weq{10}{\pm m = \overset{h,h'}{\varPi}\ \sqrt{\varkappa}\ \mathrm{sn}\,\Omega_{h,h'} = \varPi\, \varkappa_{h,h'}.}
1. \gesperrt{Die ganzen Zahlen $\varkappa_{\mu,\mu'}$ sind also Teiler der Zahl $m$.}

Wir setzen jetzt, wenn $m$ und $n$ zwei ungerade Zahlen sind:
\[
\Omega_{h,h'} = \dfrac{2hK + 2h'iK'}{m},
\]
\[
H_{l,v'} = \dfrac{2lK + 2l'iK'}{n}, % [source sic: the printed subscript is v' although the numerator uses l']
\]
und substituieren in der Formel \S\,155, (20):
\[
v = H_{l,v'}.
\]
Dadurch ergiebt sich:
\weq{11}{\dfrac{\sqrt{\varkappa}\,\mathrm{sn}\,m\,H_{l,v'}}{\sqrt{\varkappa}\,\mathrm{sn}\,H_{l,v'}} = \overset{h,h'}{\varPi}\ \sqrt{\varkappa}\ \mathrm{sn}\,(H_{l,v'} + \Omega_{h,h'}),
}
%CONT>>
